\documentclass[12pt]{book}

\usepackage[utf8]{inputenc}
\usepackage[spanish]{babel}

\usepackage{amsmath, amsthm, amsfonts}

\usepackage{xcolor}

\usepackage{listings}
\lstset{
  literate={á}{{\'a}}1
           {é}{{\'e}}1
           {í}{{\'i}}1
           {ó}{{\'o}}1
           {ú}{{\'u}}1
           {ñ}{{\~n}}1
           {ç}{{\c{c}}}1
           {Á}{{\'A}}1
           {É}{{\'E}}1
           {Í}{{\'I}}1
           {Ó}{{\'O}}1
           {Ú}{{\'U}}1
}

\begin{document}

\chapter{Código completo}

En éste apéndice pondremos el código completo sin los comentarios que mencionamos en la sección 2.3.

\begin{lstlisting}[caption={Carga de las paqueterias para la automatización, definición de la métrica (xTensor) y la carta (xCoba)}, label={lst:automatizacion-parametros}]
In[]:= << xAct`xTensor`
In[]:= << xAct`Invar`
In[]:= Block[{Print}, << xAct`xTras`]
In[]:= Block[{Print}, << xAct`xCoba`]

In[]:= (*Definimos la variedad M*)
In[]:= DefManifold[M, 4, IndexRange[a, l]] 
In[]:= DefMetric[-1, metric[-a, -b], CD, PrintAs -> "g", 
    CurvatureRelations -> True]

In[]:= DefChart[base, M, {0, 1, 2, 3}, {t[], r[], \[Theta][], \[Phi][]}, 
    ChartColor -> Red];    
\end{lstlisting}


\begin{lstlisting}[caption={}]
In[]:= DefScalarFunction[F];
In[]:= DefScalarFunction[A];
In[]:= DefScalarFunction[Q];
In[]:= DefScalarFunction[J];

In[]:= PrintAs[A] ^= "a";
In[]:= PrintAs[Q] ^= "n";
In[]:= PrintAs[J] ^= "m";
In[]:= PrintAs[F] ^= "f";

In[]:= DefConstantSymbol[\[CapitalLambda]]
In[]:= DefConstantSymbol[\[Beta]]

In[]:= PrintAs[metric] ^= "g";
In[]:= PrintAs[epsilonmetric] ^= "\     [Epsilon]";
In[]:= PrintAs[RiemannCD] ^= "R";
In[]:= PrintAs[RicciCD] ^= "R";
In[]:= PrintAs[RicciScalarCD] ^= "R";
In[]:= PrintAs[WeylCD] ^= "W";
In[]:= PrintAs[TFRicciCD] ^= "S";
\end{lstlisting}



\begin{lstlisting}[caption={Densidad lagrangiana y su variación respecto al métrica inversa}, label={lst:automatizacion-varlagrangian}]
In[]:= LagrangianDensity = RicciScalarCD[];

In[]:= EOMFirstEd = 
     VarL[metric[a, b]][LagrangianDensity] // ContractMetric // 
         FullSimplification[] // CollectTensors // ScreenDollarIndices;
\end{lstlisting}


\begin{lstlisting}[caption={Métrica para xCoba}, label={lst:automatizacion-metrica-xcoba}]
In[]:= eleccion = 
    Input["Selecciona la Métrica:\n1 - FLRW (Cosmología)\n2 - Agujero \
    Negro \n3 - Estrella "];

     G = Switch[eleccion,(*OPCIÓN 1:FLRW *)1, 
    Print["Has elegido: 1 - FLRW"];
    CTensor[
        DiagonalMatrix[{-1, (A[
            t[]])^2, (A[t[]])^2 (r[])^2, (A[
            t[]])^2 (r[])^2 (Sin[\[Theta][]])^2}], {-base, \
    -base}],(*OPCIÓN 2:AGUJERO NEGRO *)2, 
    Print["Has elegido: 2 - Agujero Negro "];
    CTensor[
        DiagonalMatrix[{-Q[r[]], Q[r[]]^-1, r[]^2, 
        r[]^2 Sin[\[Theta][]]^2}], {-base, -base}],(*OPCIÓN 3:ESTRELLA *)
    3, Print["Has elegido: 3 - Estrella "];
    CTensor[
        DiagonalMatrix[{-Q[r[]], J[r[]], r[]^2, 
        r[]^2 Sin[\[Theta][]]^2}], {-base, -base}],(*Caso por defecto \
    si escribes mal*)_, Abort[]];
\end{lstlisting}



\begin{lstlisting}[caption={Cambiando la métrica para xCoba}]
In[]:= UnsetCMetric[metric];
In[]:= SetCMetric[G, base, SignatureOfMetric -> {3, 1, 0}]; 
In[]:= MetricCompute[G, base, All];
In[]:= DC = LC[G]; 
\end{lstlisting}


\begin{lstlisting}[caption={Cambiando los tensores de xTensor a xCoba}]
In[]:= EOMThirdEd = {EOMFirstEd /. {metric -> G, 
     RicciScalarCD -> RicciScalar[DC], RicciCD -> Ricci[DC], 
     RiemannCD -> Riemann[DC]}}; 
\end{lstlisting}


\begin{lstlisting}[caption={Cambiando del lenguaje de xCoba a Mathematica nativo}]
In[]:= lenghtEOM = Length[EOMThirdEd];
In[]:= Do[
    variableEOM1 = ToBasis[base][FromBasisExpand[EOMThirdEd[[p]]]];
    variableEOM2 = TraceBasisDummy /@ variableEOM1;
     variableEOM3 = ComponentArray@variableEOM2;
    variablesEOMseparated[p] = ToValues@variableEOM3, {p, 1, lenghtEOM}
    ];
\end{lstlisting}
 

\begin{lstlisting}[caption={EOM $\mathcal{E}_{ab}$}]
In[]:= eomdd = Simplify[Sum[variablesEOMseparated[p], {p, 1, lenghtEOM}]]
\end{lstlisting}

\end{document}
